\slajdid{02-s021-rtl01}
\begin{frame}[label=02-s021-rtl01]{RTL: mreža I–ILI}
\setlength{\parskip}{0pt}
\ifdefined\DErtlDefinisan\else
\global\let\DErtlDefinisan\relax
\lstdefinelanguage{DEsystemverilog}[]{Verilog}{
 morekeywords={logic,always_comb,always_ff,timeunit,timeprecision}}
\lstdefinestyle{DErtl}{
 language=DEsystemverilog,
 basicstyle=\ttfamily\fontsize{9}{10}\selectfont,
 keywordstyle=\color{DEisticanje},
 commentstyle=\color{Muted}\itshape,
 stringstyle=\color{Ink},
 numbers=left,numberstyle=\ttfamily\fontsize{8}{10}\selectfont\color{Muted},
 numbersep=3pt,xleftmargin=10pt,
 columns=fullflexible,keepspaces=true,showstringspaces=false,
 tabsize=4,breaklines=false,numberblanklines=true,
 aboveskip=0pt,belowskip=0pt}
\fi

\begin{columns}[T,onlytextwidth]
\begin{column}{.55\textwidth}
% element: 02-s021-rtl01:d01
\lstinputlisting[style=DErtl,firstnumber=1]{kodovi/primeri_sv/02-s021/mreza_zp.sv}
\end{column}
\begin{column}{.42\textwidth}
\fontsize{10}{12}\selectfont
\textbf{Zbir proizvoda}\par\vspace{2mm}
Tri I gejta formiraju proizvode \(P1\), \(P2\) i \(P3\). Završni ILI gejt daje njihov zbir.\par\vspace{2mm}
Dva invertora po ulazu daju prave i komplementne vodove, kao u šemi.\par\vspace{2mm}
\texttt{\_n}: komplementni signal\par
\texttt{\_p}: pravi signal
\par\vspace{2mm}
% element: 02-s021-rtl01:d02
{\fontsize{9}{11}\selectfont
Testbench:\par
\href{https://github.com/pznikola/digitalna_elektronika_1/blob/main/predavanja/02_sinteza_kombinacionih_mreza/kodovi/primeri_sv/02-s021/tb_mreza_zp.sv}{\texttt{tb\_mreza\_zp.sv}}\par
Testbench bez provere:\par
\href{https://github.com/pznikola/digitalna_elektronika_1/blob/main/predavanja/02_sinteza_kombinacionih_mreza/kodovi/primeri_sv/02-s021/tb_mreza_zp_bez_provere.sv}{\texttt{tb\_mreza\_zp\_bez\_provere.sv}}\par}
\end{column}
\end{columns}
\end{frame}
